% Cyclotomic units and the plus class number.

This chapter records the cyclotomic-unit part of the proof.  The purpose is
to connect Bernoulli non-divisibility with the plus class number. Throughout
this chapter $p$ is an odd prime, $K$ is a field isomorphic to
$\Q(\zeta_p)$, and
\[
  K^+ = \Q(\zeta_p)^+
\]
is its maximal real subfield.  The group of real units is
\[
  E^+ := (\mathcal O_{K^+})^\times .
\]
Thus \(E^+\) is the full unit group of the maximal real subfield.

\section{Real cyclotomic units}

Let
\[
  u_a = \frac{1-\zeta_p^a}{1-\zeta_p}
  \qquad (a \in \Z,\ (a,p)=1)
\]
be the usual cyclotomic unit in $\mathcal O_K^\times$.  For the real
subfield one uses the conjugation-fixed unit
\[
  \varepsilon_a = u_a\,\overline{u_a}
  =
  \zeta_p^{\,1-a}
  \left(\frac{1-\zeta_p^a}{1-\zeta_p}\right)^2 .
\]
It lies in $K^+$ because complex conjugation fixes it, and it is a unit
because both $u_a$ and $\overline{u_a}$ are units.  The range used for prime
conductor is
\[
  2 \le a \le \frac{p-1}{2}.
\]
In that range $0<a<p$, so a prime $p$ cannot divide $a$; hence $(a,p)=1$ and
the above cyclotomic unit is defined.

\begin{lemma}[Coprimality of the standard indices]
  \label{lem:real-cyclotomic-index-coprime}
  \lean{BernoulliRegular.realCyclotomicUnit_index_coprime}
  \leanok
  \uses{def:galois-cyclotomic}
Let $2 \le a \le (p-1)/2$. Then $a$ is coprime to $p$.
\end{lemma}

\begin{proof}\leanok
The lower bound gives $a>0$.  The upper bound gives
$a \le (p-1)/2 < p$, so $a<p$.  If $p$ divided $a$, then since $a$ is
positive we would have $p \le a$, contradicting $a<p$.  Because $p$ is prime,
not being divisible by $p$ is equivalent to being coprime to $p$.
\end{proof}

\begin{definition}[Real cyclotomic unit]
  \label{def:real-cyclotomic-unit}
  \lean{BernoulliRegular.realCyclotomicUnit}
  \leanok
  \uses{def:cm-real-subfield, lem:real-cyclotomic-index-coprime}
For $2 \le a \le (p-1)/2$, the unit
\[
  \varepsilon_a \in E^+
\]
is the unique real unit whose image in $\mathcal O_K^\times$ is the
conjugation-fixed product $u_a\overline{u_a}$.
\end{definition}

After extension of scalars from \(\mathcal O_{K^+}\) to \(\mathcal O_K\),
this unit is exactly the standard conjugation-fixed product
\(u_a\overline{u_a}\), and is fixed by complex conjugation.

\section{The cyclotomic-unit subgroups}

There are two closely related real cyclotomic-unit subgroups.  The logarithmic
saturation argument is most convenient for the squared
family $\varepsilon_a$.  The index theorem is naturally stated for the
normalised family
\[
  \eta_a
  =
  \zeta_p^{(1-a)/2}\frac{1-\zeta_p^a}{1-\zeta_p},
  \qquad
  \eta_a^2=\varepsilon_a .
\]
Here $(1-a)/2$ is understood modulo $p$, using that $2$ is invertible modulo
the odd prime $p$.

\begin{definition}[The squared cyclotomic-unit subgroup]
  \label{def:CPlus-squared}
  \lean{BernoulliRegular.CPlus}
  \leanok
  \uses{def:real-cyclotomic-unit}
Let $g=(p-3)/2$.  The subgroup $C^+_{\mathrm{sq}}\le E^+$ is generated by
$-1$ and by the units
\[
  \varepsilon_2,\varepsilon_3,\ldots,\varepsilon_{(p-1)/2}.
\]
Equivalently,
\[
  C^+_{\mathrm{sq}}
  =
  \left\langle
    -1,\ \varepsilon_a\ (2\le a\le (p-1)/2)
  \right\rangle .
\]
\end{definition}

The finite index set may be written as
\[
  i\in \operatorname{Fin}((p-3)/2), \qquad a=i+2.
\]
Thus the listed units are precisely the nontrivial real cyclotomic generators
in the prime-conductor range, together with the torsion unit \(-1\).

\begin{definition}[The normalised cyclotomic-unit subgroup]
  \label{def:CPlus-normalized}
  \lean{BernoulliRegular.normalizedCPlus}
  \leanok
  \uses{def:real-cyclotomic-unit}
The subgroup $C^+_{\mathrm{norm}}\le E^+$ is generated by $-1$ and by the
normalised real cyclotomic units
\[
  \eta_2,\eta_3,\ldots,\eta_{(p-1)/2}.
\]
\end{definition}

The relation between the two families is exactly the relation
$\eta_a^2=\varepsilon_a$.

\begin{lemma}[Squares of normalised generators]
  \label{lem:normalized-square}
  \lean{BernoulliRegular.normalizedCPlusGenerator_sq_eq_CPlusGenerator}
  \leanok
  \uses{def:CPlus-squared, def:CPlus-normalized}
For each standard generator index $a$,
\[
  \eta_a^2=\varepsilon_a.
\]
\end{lemma}

\begin{proof}\leanok
This is the elementary identity
\[
  \left(
    \zeta_p^{(1-a)/2}\frac{1-\zeta_p^a}{1-\zeta_p}
  \right)^2
  =
  \zeta_p^{1-a}
  \left(\frac{1-\zeta_p^a}{1-\zeta_p}\right)^2
  =
  u_a\overline{u_a}.
\]
The last equality is the standard computation
\[
  \overline{u_a}
  =
  \frac{1-\zeta_p^{-a}}{1-\zeta_p^{-1}}
  =
  \zeta_p^{1-a}\frac{1-\zeta_p^a}{1-\zeta_p}.
\]
\end{proof}

\begin{lemma}[Odd-primary comparison of the two indices]
  \label{lem:CPlus-index-normalized-index}
  \lean{BernoulliRegular.CPlus_index_prime_dvd_iff_normalizedCPlus_index_prime_dvd}
  \lean{BernoulliRegular.subgroup_index_prime_dvd_iff_of_relIndex_dvd_two_pow}
  \leanok
  \uses{def:CPlus-squared, def:CPlus-normalized, lem:normalized-square}
For odd $p$,
\[
  p \mid [E^+ : C^+_{\mathrm{sq}}]
  \quad\Longleftrightarrow\quad
  p \mid [E^+ : C^+_{\mathrm{norm}}].
\]
\end{lemma}

\begin{proof}\leanok
The subgroup $C^+_{\mathrm{sq}}$ is contained in
$C^+_{\mathrm{norm}}$, because $-1$ is a normalised generator and
each squared generator $\varepsilon_a$ is the square of the normalised
generator $\eta_a$.  Conversely, every class in
$C^+_{\mathrm{norm}}/C^+_{\mathrm{sq}}$ has square one: it is enough to check
this on the generators, and it follows from $\eta_a^2=\varepsilon_a$ and
$(-1)^2=1$.  Hence the relative index
$[C^+_{\mathrm{norm}}:C^+_{\mathrm{sq}}]$ divides a power of $2$.

The two ambient indices differ by this relative index.  Since $p$ is odd, $p$
does not divide any power of $2$.  Therefore the $p$-primary divisibility of
$[E^+ : C^+_{\mathrm{sq}}]$ and $[E^+ : C^+_{\mathrm{norm}}]$ is the same.
\end{proof}

\section{Exact p-saturation}

Let $G$ be an abelian group and $H\le E\le G$.  We use the exact image of the
$p$-th power map:
\[
  H^p := \{h^p : h\in H\}.
\]
This is not a closure or saturation of the image; it is the literal subgroup
of $p$-th powers.

\begin{definition}[Exact $p$-saturation]
  \label{def:p-saturated}
  \lean{BernoulliRegular.pSaturated}
  \lean{BernoulliRegular.pPowerSubgroup}
  \leanok
The subgroup $H$ is $p$-saturated in $E$ if $H\le E$ and
\[
  H\cap E^p \subseteq H^p.
\]
Equivalently, whenever $x\in H$ and $x=y^p$ for some $y\in E$, there exists
$z\in H$ such that $x=z^p$.
\end{definition}

The useful criterion for $C^+_{\mathrm{sq}}$ is stated in terms of exponent
vectors.  Every element of $C^+_{\mathrm{sq}}$ can be written as
\[
  (-1)^s\prod_{a=2}^{(p-1)/2}\varepsilon_a^{e_a},
  \qquad s,e_a\in\Z .
\]
The following criterion uses these exponent vectors.

\begin{theorem}[Exponent criterion for saturation]
  \label{thm:CPlus-pSaturated-exponents}
  \lean{BernoulliRegular.exists_CPlusExponentProduct_of_mem_CPlus}
  \lean{BernoulliRegular.CPlus_pSaturated_of_generator_exponents_modP_zero}
  \lean{BernoulliRegular.CPlusExponentProduct_mem_CPlus}
  \lean{BernoulliRegular.neg_one_zpow_mem_pPowerSubgroup_CPlus}
  \lean{BernoulliRegular.CPlusExponentProduct_pow_of_exponents_eq_mul}
  \leanok
  \uses{def:p-saturated, def:CPlus-squared}
Assume that for every expression
\[
  x=(-1)^s\prod_a \varepsilon_a^{e_a}\in C^+_{\mathrm{sq}}
\]
which is a $p$-th power in $E^+$, all exponents satisfy
\[
  e_a\equiv 0 \pmod p.
\]
Then $C^+_{\mathrm{sq}}$ is $p$-saturated in $E^+$.
\end{theorem}

\begin{proof}\leanok
Let $x\in C^+_{\mathrm{sq}}$ and suppose $x=y^p$ for some $y\in E^+$.
Choose an exponent expression
\[
  x=(-1)^s\prod_a \varepsilon_a^{e_a}.
\]
By hypothesis, $e_a\equiv 0\pmod p$ for every $a$, so choose integers
$k_a$ with $e_a=pk_a$.  Since $p$ is odd, the sign term is already a
$p$-th power of itself:
\[
  ((-1)^s)^p=(-1)^s.
\]
Therefore
\[
  x
  =
  (-1)^s\prod_a \varepsilon_a^{pk_a}
  =
  \left((-1)^s\prod_a \varepsilon_a^{k_a}\right)^p .
\]
The element inside the parentheses belongs to $C^+_{\mathrm{sq}}$, so $x$ is
a $p$-th power inside $C^+_{\mathrm{sq}}$.  This is exactly
$p$-saturation.
\end{proof}

\section{From saturation to index non-divisibility}

\begin{theorem}[Saturation forbids $p$ in the index]
  \label{thm:not-dvd-index-of-pSaturated}
  \lean{BernoulliRegular.not_dvd_index_of_pSaturated}
  \lean{BernoulliRegular.subgroup_not_dvd_index_of_pSaturated_top_of_pow_eq_one_mem}
  \lean{BernoulliRegular.CPlus_index_ne_zero}
  \leanok
  \uses{def:p-saturated, def:CPlus-squared}
If $C^+_{\mathrm{sq}}$ is $p$-saturated in $E^+$, then
\[
  p\nmid [E^+:C^+_{\mathrm{sq}}].
\]
\end{theorem}

\begin{proof}\leanok
The proof is purely group-theoretic once one knows that
$C^+_{\mathrm{sq}}$ has finite index in $E^+$ and contains the torsion of
$E^+$.  Suppose for contradiction that $p$ divides
$[E^+:C^+_{\mathrm{sq}}]$.  By Cauchy's theorem applied to the finite quotient
$E^+/C^+_{\mathrm{sq}}$, there is a non-trivial quotient class of order $p$.
Choose a representative $g\in E^+$.  The order condition says
\[
  g^p\in C^+_{\mathrm{sq}}.
\]
Of course $g^p$ is a $p$-th power in $E^+$.  By $p$-saturation, there is
$h\in C^+_{\mathrm{sq}}$ with
\[
  h^p=g^p.
\]
Then
\[
  (gh^{-1})^p=1.
\]
Thus $gh^{-1}$ is $p$-torsion in $E^+$.  The only torsion units in the real
cyclotomic field are $\pm1$, and both belong to $C^+_{\mathrm{sq}}$.  Hence
$gh^{-1}\in C^+_{\mathrm{sq}}$, and since $h\in C^+_{\mathrm{sq}}$, also
$g\in C^+_{\mathrm{sq}}$.  This means the quotient class of $g$ was trivial,
contradicting its order $p$.
\end{proof}

\section{The Kummer logarithm determinant}

The $p$-adic part of the argument studies the logarithms of the real
cyclotomic units at the prime above $p$.  After passing to the completed local
ring and reducing suitable logarithmic coordinates modulo $p$, one obtains a
matrix over $\mathbb F_p$.  Its columns are indexed by the generators
\[
  \varepsilon_2,\ldots,\varepsilon_{(p-1)/2},
\]
and its rows are indexed by the even powers which occur in the Kummer range.

The determinant has two parts.  One part is a finite-field Vandermonde
determinant in the Teichmuller nodes; this is nonzero because the nodes are
distinct.  The other part is the product of Bernoulli factors.  Consequently
the determinant is nonzero exactly when no Bernoulli numerator in Kummer's
range is divisible by $p$.

\begin{theorem}[Kummer logarithm determinant criterion]
  \label{thm:kummer-log-det-bernoulli}
  \lean{BernoulliRegular.CyclotomicUnits.kummerLogMatrix_det_ne_zero_iff_bernoulli_nonzero}
  \lean{BernoulliRegular.CyclotomicUnits.concreteKummerLogMatrix_eq_diagonal_mul_vandermonde}
  \lean{BernoulliRegular.CyclotomicUnits.concreteKummerLogMatrix_det_ne_zero_iff_forall_bernoulliFactor_ne_zero}
  \lean{BernoulliRegular.CyclotomicUnits.bernoulliFactor_ne_zero_iff_not_dvd_bernoulli_num}
  \lean{BernoulliRegular.CyclotomicUnits.vandermonde_teichmuller_even_sub_one_det_ne_zero}
  \leanok
  \uses{def:CPlus-squared, thm:bernoulli-padicInt-below, cor:bernoulli-den-coprime}
Assume $p\ge 5$.  The Kummer logarithm matrix has nonzero determinant if and
only if
\[
  \forall j,\quad
  1\le j,\ 2j\le p-3
  \Longrightarrow
  p\nmid (B_{2j})_{\mathrm{num}}.
\]
\end{theorem}

\begin{proof}\leanok
The determinant computation factors the matrix as
\[
  \operatorname{diag}(r_1,\ldots,r_g)\,V,
  \qquad g=(p-3)/2,
\]
where $V$ is the finite-field Vandermonde matrix attached to the even
Teichmuller nodes and $r_j$ is the Bernoulli row factor attached to
$B_{2j}/(2j)$.  The determinant is therefore
\[
  \det(V)\prod_j r_j.
\]
The Vandermonde determinant is nonzero because the Teichmuller nodes are
pairwise distinct in $\mathbb F_p$.  Thus the determinant is nonzero exactly
when every $r_j$ is nonzero.

For $1\le j$ and $2j\le p-3$, the denominator of $B_{2j}$ is prime to $p$ by
the integrality results for Bernoulli numbers below the boundary.  Also
$2j<p$.  Hence reducing $B_{2j}/(2j)$ modulo $p$ gives zero exactly when the
integer numerator of $B_{2j}$ is divisible by $p$.  This identifies the
non-vanishing of all row factors with the displayed Bernoulli
non-divisibility condition.
\end{proof}

\begin{theorem}[Determinant non-vanishing gives saturation]
  \label{thm:kummer-log-det-saturation}
  \lean{BernoulliRegular.CyclotomicUnits.cyclotomicUnits_pSaturated_of_kummerLog_det_ne_zero}
  \lean{BernoulliRegular.CyclotomicUnits.completedLog_relation_of_CPlus_product_mem_powers}
  \lean{BernoulliRegular.CyclotomicUnits.concreteKummerLogMatrix_mulVec_exponents_eq_zero}
  \lean{BernoulliRegular.CyclotomicUnits.exponents_modP_eq_zero_of_kummerLogMatrix_relation}
  \lean{BernoulliRegular.CyclotomicUnits.CPlusGenerator_exponents_modP_zero_of_kummerLog_det_ne_zero}
  \leanok
  \uses{def:p-saturated, def:CPlus-squared, thm:CPlus-pSaturated-exponents}
Assume $p\ge 5$.  If the Kummer logarithm determinant is nonzero, then
$C^+_{\mathrm{sq}}$ is $p$-saturated in $E^+$.
\end{theorem}

\begin{proof}\leanok
Suppose an element
\[
  x=(-1)^s\prod_a \varepsilon_a^{e_a}
\]
is a $p$-th power in $E^+$.  Taking completed $p$-adic logarithms of the
local image of this relation gives an additive relation
\[
  \sum_a e_a\,\log(\varepsilon_a)
  \in p\cdot L,
\]
where $L$ is the completed logarithm lattice. The
sign term disappears from the logarithmic coordinates, because $p$ is odd and
the sign is torsion.

Reducing the coordinates of this relation modulo $p$ gives
\[
  M\,(e_a\bmod p)_a=0
\]
for the Kummer logarithm matrix $M$.  If $\det(M)\ne0$, then $M$ is invertible
over $\mathbb F_p$, so every coordinate $e_a\bmod p$ is zero.  The exponent
criterion of \cref{thm:CPlus-pSaturated-exponents} then proves
$p$-saturation.
\end{proof}

\section{Bernoulli non-divisibility and the plus index}

The previous sections combine into the theorem used by the final proof.

\begin{theorem}[Bernoulli non-divisibility implies index non-divisibility]
  \label{thm:bernoulli-nonzero-index}
  \lean{BernoulliRegular.not_dvd_cyclotomicUnitIndex_of_bernoulli_nonzero}
  \leanok
  \uses{thm:kummer-log-det-bernoulli, thm:kummer-log-det-saturation,
    thm:not-dvd-index-of-pSaturated}
Assume $p$ is odd and
\[
  \forall j,\quad
  1\le j,\ 2j\le p-3
  \Longrightarrow
  p\nmid (B_{2j})_{\mathrm{num}}.
\]
Then
\[
  p\nmid [E^+:C^+_{\mathrm{sq}}].
\]
\end{theorem}

\begin{proof}\leanok
If $p=3$, then the generator range is empty.  The real subfield of
$\Q(\zeta_3)$ is $\Q$, so its unit group is generated by $-1$.  Hence
$C^+_{\mathrm{sq}}=E^+$ and the index is $1$, which is not divisible by $p$.

Now assume $p\ge5$.  The Bernoulli non-divisibility hypothesis makes the
Kummer logarithm determinant nonzero by
\cref{thm:kummer-log-det-bernoulli}.  The determinant non-vanishing theorem
\cref{thm:kummer-log-det-saturation} gives $p$-saturation of
$C^+_{\mathrm{sq}}$ in $E^+$.  Finally
\cref{thm:not-dvd-index-of-pSaturated} converts this saturation into
\[
  p\nmid [E^+:C^+_{\mathrm{sq}}].
\]
\end{proof}

\section{The needed Sinnott index formula}

The index theorem used in the final argument is the odd-primary consequence
of Sinnott's formula for real cyclotomic units.  The family used in the
regulator computation is the same squared family as above, but with the
torsion subgroup adjoined explicitly.

\begin{definition}[Sinnott's full-rank subgroup]
  \label{def:sinnott-index-subgroup}
  \lean{BernoulliRegular.cyclotomicUnitIndexSubgroup}
  \leanok
  \uses{def:real-cyclotomic-unit, def:CPlus-squared}
For $p\ge3$, let
\[
  C^+_{\mathrm{Sin}}
  =
  \left\langle
    \varepsilon_2,\varepsilon_3,\ldots,\varepsilon_{(p-1)/2}
  \right\rangle \vee E^+_{\mathrm{tors}}
  \le E^+ .
\]
Equivalently, $C^+_{\mathrm{Sin}}$ is generated by the full-rank real
cyclotomic-unit family used in the regulator determinant, together with all
torsion units of $E^+$.
\end{definition}

\begin{lemma}[Agreement with the squared subgroup]
  \label{lem:sinnott-subgroup-eq-CPlus}
  \lean{BernoulliRegular.cyclotomicUnitIndexSubgroup_eq_CPlus}
  \lean{BernoulliRegular.CPlusGenerator_eq_cyclotomicUnitFamilyKplus}
  \lean{BernoulliRegular.range_cyclotomicUnitFamilyKplusFinRank_eq}
  \lean{BernoulliRegular.torsionKplus_le_CPlus}
  \lean{BernoulliRegular.CPlus_le_cyclotomicUnitIndexSubgroup}
  \lean{BernoulliRegular.cyclotomicUnitIndexSubgroup_le_CPlus}
  \leanok
  \uses{def:sinnott-index-subgroup, def:CPlus-squared}
For $p\ge3$,
\[
  C^+_{\mathrm{Sin}}=C^+_{\mathrm{sq}}.
\]
\end{lemma}

\begin{proof}\leanok
The full-rank family in $C^+_{\mathrm{Sin}}$ is indexed by
$\operatorname{Fin}((p-3)/2)$, and the element with index $i$ is exactly
$\varepsilon_{i+2}$.  Thus its range is the same generator range used for
$C^+_{\mathrm{sq}}$.

It remains to compare torsion.  The torsion of the real cyclotomic unit group
is $\{\pm1\}$.  The unit $1$ belongs to every subgroup, and $-1$ is one of the
generators of $C^+_{\mathrm{sq}}$.  Hence adjoining torsion to the full-rank
family adds nothing beyond the subgroup already generated by
$-1,\varepsilon_2,\ldots,\varepsilon_{(p-1)/2}$.
\end{proof}

\begin{theorem}[Sinnott formula, odd-primary consequence]
  \label{thm:sinnott-index-pprimary}
  \lean{BernoulliRegular.cyclotomicUnitIndex_primeConductor_pPrimary_of_sinnottIndexFormula}
  \lean{BernoulliRegular.cyclotomicUnitIndex_primeConductor_pPrimary_of_kummerDirichletDeterminant}
  \lean{BernoulliRegular.FLT37.Sinnott.sinnottRegulatorIdentity_iff_kummerDirichletDeterminant}
  \lean{BernoulliRegular.FLT37.Sinnott.sinnottIndexFormula_of_regulatorIdentity}
  \lean{BernoulliRegular.FLT37.Sinnott.index_eq_twoPow_mul_hPlus_of_sinnottIndexFormula}
  \lean{BernoulliRegular.prime_dvd_two_pow_mul_iff_dvd}
  \leanok
  \uses{def:sinnott-index-subgroup}
Assume the Kummer-Dirichlet determinant identity for the real
cyclotomic-unit logarithm matrix.  Then
\[
  p\mid [E^+:C^+_{\mathrm{Sin}}]
  \quad\Longleftrightarrow\quad
  p\mid h^+ .
\]
\end{theorem}

\begin{proof}\leanok
Write $g=(p-3)/2$ and let
\[
  R^+ = \operatorname{Reg}(E^+)
\]
be the regulator of the maximal real subfield.  The determinant identity is
the regulator identity
\[
  \operatorname{Reg}(\varepsilon_2,\ldots,\varepsilon_{(p-1)/2})
  =
  2^g h^+ R^+ .
\]

The passage from regulators to indices is the standard covolume comparison
for the logarithmic unit lattice.  Modulo torsion, the units $E^+$ form a
lattice of rank $g$, and the family
$\varepsilon_2,\ldots,\varepsilon_{(p-1)/2}$ spans a full-rank sublattice.
The covolume of the latter is its index in the former times the covolume of
the former:
\[
  \frac{
    \operatorname{Reg}(\varepsilon_2,\ldots,\varepsilon_{(p-1)/2})
  }{R^+}
  =
  [E^+:C^+_{\mathrm{Sin}}].
\]
Since $R^+>0$, the displayed regulator identity therefore gives Sinnott's
index formula in this normalization:
\[
  [E^+:C^+_{\mathrm{Sin}}]=2^g h^+ .
\]

Finally, $p$ is odd, so $p\nmid 2^g$.  Euclid's lemma gives
\[
  p\mid 2^g h^+
  \quad\Longleftrightarrow\quad
  p\mid h^+,
\]
and combining this with the index formula proves the stated equivalence.
\end{proof}

\begin{theorem}[Deleted-Fourier determinant identity]
  \label{thm:deleted-fourier-kummer-dirichlet}
  \lean{BernoulliRegular.kummerDirichletDeterminant_of_deletedFourier}
  \leanok
  \uses{def:sinnott-index-subgroup, def:dirichlet-character,
    def:real-cyclotomic-unit}
Assume $p\ge5$.  The logarithm matrix of the full-rank real cyclotomic-unit
family satisfies the Kummer-Dirichlet determinant identity.
\end{theorem}

\begin{proof}\leanok
The rows of the matrix are indexed by the non-trivial even characters modulo
$p$, and the columns by
\[
  \varepsilon_2,\ldots,\varepsilon_{(p-1)/2}.
\]
After the standard Fourier change of basis, the determinant separates into a
deleted Fourier determinant and the product of the even character factors.

The deleted Fourier determinant is nonzero: it is the determinant of the
character table with the trivial row and the redundant column removed, and
orthogonality of characters gives its inverse after multiplying by a power of
$p$.  The remaining product is the even part of the cyclotomic class-number
formula.  With the present squared-unit normalization, the regulator identity
therefore has the form
\[
  \operatorname{Reg}(\varepsilon_2,\ldots,\varepsilon_{(p-1)/2})
  =
  2^g h^+ R^+ ,
  \qquad g=(p-3)/2.
\]
This is precisely the Kummer-Dirichlet determinant identity used in the
index calculation.
\end{proof}

\begin{corollary}[Sinnott formula for prime conductor]
  \label{cor:sinnott-index-prime-conductor}
  \lean{BernoulliRegular.kummerDirichletDeterminant_of_deletedFourier}
  \lean{BernoulliRegular.cyclotomicUnitIndex_primeConductor_pPrimary_aux}
  \lean{BernoulliRegular.cyclotomicUnitIndex_primeConductor_pPrimary_of_kummerDirichletDeterminant}
  \leanok
  \uses{thm:sinnott-index-pprimary, thm:deleted-fourier-kummer-dirichlet}
If $p\ge5$, then
\[
  p\mid [E^+:C^+_{\mathrm{Sin}}]
  \quad\Longleftrightarrow\quad
  p\mid h^+ .
\]
\end{corollary}

\begin{proof}\leanok
For $p\ge5$, the deleted-Fourier determinant computation evaluates the
logarithm determinant of the real cyclotomic units.  The matrix rows are the
even non-trivial characters, the columns are the units
$\varepsilon_2,\ldots,\varepsilon_{(p-1)/2}$, and the determinant splits into
a nonzero Fourier determinant times the product of the character factors.
The class-number formula for the real subfield identifies this product with
$h^+R^+$, with the extra factor $2^g$ coming from using the squared
cyclotomic-unit family.  Thus the Kummer-Dirichlet determinant identity
holds, and \cref{thm:sinnott-index-pprimary} applies.
\end{proof}

\begin{theorem}[Prime-conductor cyclotomic-unit index theorem]
  \label{thm:cyclotomic-unit-index-hplus}
  \lean{BernoulliRegular.cyclotomicUnitIndex_primeConductor_pPrimary}
  \lean{BernoulliRegular.cyclotomicUnitIndex_primeConductor_pPrimary_of_eq_three}
  \lean{BernoulliRegular.cyclotomicUnitIndex_primeConductor_pPrimary_of_five_le}
  \leanok
  \uses{cor:sinnott-index-prime-conductor, lem:sinnott-subgroup-eq-CPlus,
    lem:CPlus-index-normalized-index}
For odd $p$,
\[
  p \mid [E^+:C^+_{\mathrm{norm}}]
  \quad\Longleftrightarrow\quad
  p \mid h^+.
\]
\end{theorem}

\begin{proof}\leanok
There are two cases.

If \(p=3\), then \(K^+=\Q\).  The plus class number is therefore \(h^+=1\).
The real unit rank is zero, so the normalised cyclotomic-unit subgroup is the
whole unit group.  Hence both sides are the assertion that \(3\) divides \(1\),
and both are false.

Now assume $p\ne3$.  Since $p$ is an odd prime, this gives $p\ge5$.  By
\cref{cor:sinnott-index-prime-conductor},
\[
  p\mid [E^+:C^+_{\mathrm{Sin}}]
  \quad\Longleftrightarrow\quad
  p\mid h^+ .
\]
The subgroup identification
\cref{lem:sinnott-subgroup-eq-CPlus} replaces $C^+_{\mathrm{Sin}}$ by
$C^+_{\mathrm{sq}}$.  Finally, the odd-primary comparison
\cref{lem:CPlus-index-normalized-index} replaces
$C^+_{\mathrm{sq}}$ by $C^+_{\mathrm{norm}}$, because the quotient between
the two generated subgroups has exponent dividing a power of $2$.  Since
$p$ is odd, this quotient does not affect $p$-divisibility of the ambient
index.  The displayed equivalence follows.
\end{proof}
